% =====================================================================
%  Document de BRAINSTORMING - Architecture du systeme d'analyse de
%  sentiments (Projet d'application en Data Mining - ENEAM/UAC)
%  Compilation : tectonic brainstorming_architecture.tex
%  Police : Arial (rendue via Liberation Sans, metriquement identique)
% =====================================================================
\documentclass[11pt,a4paper]{article}

\usepackage{fontspec}
\setmainfont{Liberation Sans}   % equivalent metrique exact d'Arial
\setmonofont{Liberation Mono}
\usepackage[french]{babel}

\usepackage{geometry}
\geometry{margin=2.2cm}
\usepackage{xcolor}
\usepackage{graphicx}
\usepackage{booktabs}
\usepackage{tabularx}
\usepackage{array}
\usepackage{enumitem}
\usepackage{titlesec}
\usepackage{fancyhdr}
\usepackage[most]{tcolorbox}
\tcbset{after={\par\normalcolor\medskip}}   % securite : evite la fuite de couleur (xetex)
\usepackage{hyperref}

% --- Couleurs ---
\definecolor{bleu}{RGB}{20,70,130}
\definecolor{bleuclair}{RGB}{232,240,250}
\definecolor{vert}{RGB}{30,120,70}
\definecolor{vertclair}{RGB}{232,246,238}
\definecolor{orange}{RGB}{200,110,20}
\definecolor{orangeclair}{RGB}{253,243,230}
\definecolor{gris}{RGB}{90,90,90}

\hypersetup{colorlinks=true, linkcolor=bleu, urlcolor=bleu, citecolor=bleu}

% --- Titres ---
\titleformat{\section}{\Large\bfseries\color{bleu}}{\thesection.}{0.5em}{}
\titleformat{\subsection}{\large\bfseries\color{bleu!80!black}}{\thesubsection}{0.5em}{}
\titlespacing*{\section}{0pt}{1.4em}{0.6em}

% --- En-tete / pied ---
\pagestyle{fancy}
\fancyhf{}
\fancyhead[L]{\footnotesize\color{gris}\itshape Analyse de sentiments -- Brainstorming}
\fancyhead[R]{\footnotesize\color{gris}ENEAM/UAC -- M2 ID}
\fancyfoot[C]{\thepage}
\renewcommand{\headrulewidth}{0.4pt}

% --- Boites ---
\newtcolorbox{decision}{colback=bleuclair,colframe=bleu,boxrule=0.6pt,arc=2mm,
  left=2mm,right=2mm,top=1mm,bottom=1mm,
  fonttitle=\bfseries,title=Décision de conception}
\newtcolorbox{repere}{colback=vertclair,colframe=vert,boxrule=0.6pt,arc=2mm,
  left=2mm,right=2mm,top=1mm,bottom=1mm,
  fonttitle=\bfseries,title=Exigence du cahier des charges}
\newtcolorbox{alerte}{colback=orangeclair,colframe=orange,boxrule=0.6pt,arc=2mm,
  left=2mm,right=2mm,top=1mm,bottom=1mm,
  fonttitle=\bfseries,title=Point de vigilance}

\newcommand{\code}[1]{\texttt{\small #1}}

% =====================================================================
\begin{document}

\begin{titlepage}
\centering
\vspace*{2cm}
{\color{bleu}\rule{\linewidth}{1.2pt}}\\[0.6cm]
{\Huge\bfseries Système d'analyse de sentiments\\ et d'opinions en français}\\[0.4cm]
{\Large Document de brainstorming -- Architecture complète}\\[0.3cm]
{\large\itshape Projet d'application en Data Mining}\\[0.2cm]
{\color{bleu}\rule{\linewidth}{1.2pt}}\\[1.2cm]

\begin{tabular}{rl}
\textbf{UFR} & ENEAM / UAC\\
\textbf{Niveau} & Master 2 -- Informatique Décisionnelle\\
\textbf{Année} & 2025--2026\\
\textbf{Système retenu} & Analyse de sentiments et d'opinions\\
\textbf{Langue cible} & Français\\
\end{tabular}

\vfill
{\color{gris}\small Document de travail -- conception et choix d'architecture}
\end{titlepage}

\tableofcontents
\newpage

% =====================================================================
\section{Compréhension et recadrage du sujet}

Le cahier des charges demande de \textbf{concevoir une application graphique}
(Python ou Java) pour l'un des deux systèmes de data mining proposés. Nous
retenons le \textbf{système d'analyse de sentiments et d'opinions}.

\begin{repere}
L'application doit pouvoir analyser \textbf{trois types d'entrées} :
\textbf{(1)} des \emph{phrases isolées}, \textbf{(2)} des \emph{textes / documents},
\textbf{(3)} des \emph{données en provenance des réseaux sociaux} --- le tout
\textbf{en français}. Elle doit s'appuyer sur les \textbf{deux grandes approches}
citées : l'approche \textbf{Machine Learning} et l'approche \textbf{lexicale}.
\end{repere}

C'est le point structurant : le projet n'est \emph{pas} un simple pipeline de
collecte de tweets, mais une \textbf{application interactive} dont le moteur
combine deux familles de méthodes, avec une finalité d'\textbf{aide à la
décision} (recherche marketing, supervision des réseaux sociaux, prise de
décision), exactement comme l'énumère le sujet.

\subsection{Reformulation des objectifs}
\begin{itemize}[leftmargin=1.2cm,nosep]
  \item Classer un contenu français en \textbf{positif / neutre / négatif}
        (et, en option, détecter l'\emph{émotion} et les \emph{aspects}).
  \item Offrir une \textbf{interface graphique} couvrant les 3 modes d'entrée.
  \item Implémenter et \textbf{comparer} l'approche lexicale et l'approche ML,
        puis sélectionner le modèle \textbf{optimal} sur des métriques objectives.
  \item Produire des \textbf{visualisations décisionnelles} (distribution,
        tendances, termes saillants) et un export exploitable.
\end{itemize}

% =====================================================================
\section{Vue d'ensemble de l'architecture}

L'application est organisée en couches faiblement couplées. Le \emph{moteur de
sentiment} est le c\oe ur ; l'interface et les sources ne sont que des
adaptateurs autour de lui.

{\scriptsize
\begin{verbatim}
 +---------------------------------------------------------------+
 |                 INTERFACE GRAPHIQUE (Streamlit)               |
 |   [Onglet 1] Phrase   [Onglet 2] Texte/Fichier   [Onglet 3]  |
 |                                            Reseaux sociaux    |
 +-------------------------------+-------------------------------+
                                 |
                      +----------v-----------+
                      |  COUCHE D'ENTREE     |  phrase | fichier | collecte
                      |  (adaptateurs)       |
                      +----------+-----------+
                                 |
                      +----------v-----------+
                      |  PRETRAITEMENT FR    |  nettoyage, negation,
                      |  (NLP francais)      |  emojis, tokenisation
                      +----------+-----------+
                                 |
            +--------------------+--------------------+
            |          MOTEUR DE SENTIMENT           |
            |  (A) Approche LEXICALE   (B) Approche ML|
            |   lexique FR + regles    classique +    |
            |                          transformers   |
            +--------------------+--------------------+
                                 |
                      +----------v-----------+
                      |  EVALUATION & CHOIX  |  metriques, modele optimal
                      +----------+-----------+
                                 |
                      +----------v-----------+
                      |  RESTITUTION /       |  graphes, tendances,
                      |  AIDE A LA DECISION  |  export CSV/PDF
                      +----------------------+
\end{verbatim}}

% =====================================================================
\section{Couche 1 -- Interface graphique}

\begin{decision}
\textbf{Streamlit} comme interface principale (web, Python pur, déjà présent
dans l'environnement \code{datamining}). Alternatives possibles :
\textbf{NiceGUI} (web, plus riche) ou \textbf{CustomTkinter} (bureau natif).
\end{decision}

Trois onglets, alignés sur les trois modes exigés :
\begin{enumerate}[leftmargin=1.2cm,nosep]
  \item \textbf{Phrase isolée} : champ de saisie $\rightarrow$ polarité + score
        + justification (mots déclencheurs surlignés).
  \item \textbf{Texte / document} : import \code{.txt}, \code{.csv}, \code{.pdf}
        $\rightarrow$ analyse phrase par phrase + synthèse agrégée.
  \item \textbf{Réseaux sociaux} : collecte (scraper) \emph{ou} import d'un jeu
        de données $\rightarrow$ analyse de masse, tendances, termes saillants.
\end{enumerate}
Un \textbf{sélecteur de moteur} (lexique / ML classique / transformer) permet de
basculer entre les approches et de \textbf{comparer} les résultats en direct ---
ce qui matérialise visuellement l'exigence du sujet.

% =====================================================================
\section{Couche 2 -- Prétraitement du français}

Étape critique car la qualité de l'analyse en dépend.
\begin{itemize}[leftmargin=1.2cm,nosep]
  \item Normalisation : minuscules, URLs / mentions / hashtags, espaces.
  \item \textbf{Emojis} : conversion en tokens de sentiment (porteurs forts).
  \item \textbf{Négation} : marquage (\og ne\dots pas \fg{}) car elle inverse la
        polarité --- indispensable en français.
  \item \textbf{Intensifieurs} (\og très \fg{}, \og vraiment \fg{}) et atténuateurs.
  \item Tokenisation / lemmatisation : \textbf{spaCy} (\code{fr\_core\_news\_md})
        ou NLTK.
\end{itemize}

% =====================================================================
\section{Couche 3 -- Moteur de sentiment (le c\oe ur du projet)}

Le sujet impose les \textbf{deux approches}. Nous les implémentons toutes les
deux, ce qui permet une \textbf{comparaison} et donne sa valeur scientifique au
rapport.

\subsection{(A) Approche lexicale}
\begin{itemize}[leftmargin=1.2cm,nosep]
  \item Ressource : un \textbf{lexique de sentiments français} ---
        \textbf{FEEL} (French Expanded Emotion Lexicon), \textbf{Polarimots},
        ou un portage de \textbf{VADER} pour le français.
  \item Calcul : score = somme pondérée des polarités des termes, \textbf{modulé}
        par la négation, les intensifieurs et les emojis.
  \item Atouts : \textbf{zéro entraînement}, totalement \textbf{explicable}
        (on montre les mots qui ont compté). Sert de \textbf{baseline}.
  \item Limite : faible sur l'ironie, le contexte, le vocabulaire nouveau.
\end{itemize}

\subsection{(B) Approche Machine Learning}
Deux niveaux, du plus \og data mining classique \fg{} au plus performant.

\paragraph{B.1 -- ML classique (c\oe ur \og data mining \fg{}).}
Vectorisation \textbf{TF-IDF / sac-de-mots} puis classifieurs :
\textbf{Naïve Bayes multinomial}, \textbf{SVM linéaire}, \textbf{Régression
logistique}, \textbf{Random Forest} (et \textbf{XGBoost} en option). C'est la
partie qui mobilise directement les techniques du cours (apprentissage
supervisé, évaluation, sélection de modèle).

\paragraph{B.2 -- Transformers (état de l'art français).}
\textbf{CamemBERT} (ou \textbf{DistilCamemBERT}, plus léger) \textbf{fine-tuné}
sur un corpus de sentiment français, via la librairie \textbf{Hugging Face
Transformers} (\code{Trainer}). Meilleure précision, gère mieux le contexte.

\begin{decision}
\textbf{Volet \og innovation \fg{} (facultatif mais valorisant) :} fine-tuning
d'un \textbf{petit LLM français} (type Mistral / Llama instruct) avec
\textbf{Unsloth} en \textbf{QLoRA 4-bit}. Unsloth divise par $\approx$2 le temps
d'entraînement et l'empreinte mémoire --- adapté à un GPU modeste. Intérêt :
au-delà de la classe, le LLM peut produire une \textbf{justification} et extraire
les \textbf{aspects} (\og le service est lent mais le produit est bon \fg{}).
\end{decision}

\subsection{Synthèse comparative des approches}
\begin{center}
\small
\begin{tabularx}{\linewidth}{>{\bfseries}l c c X}
\toprule
Approche & Entraînement & Explicabilité & Performance attendue (FR) \\
\midrule
Lexicale (FEEL) & Aucun & Très forte & Moyenne (baseline) \\
ML classique TF-IDF & Léger & Forte & Bonne \\
CamemBERT fine-tuné & Moyen (GPU) & Faible & Très bonne \\
LLM + Unsloth (QLoRA) & Lourd (GPU) & Moyenne (justifié) & Très bonne + aspects \\
\bottomrule
\end{tabularx}
\end{center}

% =====================================================================
\section{Couche 4 -- Données, entraînement et choix du modèle optimal}

\subsection{Jeux de données français}
\begin{itemize}[leftmargin=1.2cm,nosep]
  \item \textbf{Allociné} : critiques de films (polarité binaire) -- grand, propre.
  \item \textbf{Avis / tweets français annotés} (Hugging Face Hub : \code{tblard/allocine},
        corpus de tweets FR, etc.).
  \item \textbf{Corpus maison} : quelques centaines de messages annotés à la main
        (positif/neutre/négatif) pour évaluer sur \emph{ton} domaine réel.
\end{itemize}

\subsection{Protocole de sélection (le \og modèle optimal \fg{})}
\begin{enumerate}[leftmargin=1.2cm,nosep]
  \item Découpage \textbf{train / validation / test} + \textbf{validation croisée}.
  \item Métriques : \textbf{exactitude}, \textbf{précision}, \textbf{rappel},
        \textbf{F1-macro}, \textbf{matrice de confusion}.
  \item \textbf{Comparaison} : lexique vs ML classique vs CamemBERT (vs LLM).
  \item Le meilleur F1-macro sur le jeu de test = \textbf{modèle retenu} en
        production dans l'application.
\end{enumerate}
\begin{repere}
C'est cette démarche comparative et chiffrée qui constitue le \og data mining \fg{}
attendu : on ne choisit pas un modèle au hasard, on \textbf{démontre} son
optimalité.
\end{repere}

\subsection{Outils d'entraînement}
\code{scikit-learn} (ML classique), \code{Hugging Face Transformers}~+~\code{PyTorch}
(CamemBERT), \code{Unsloth}~+~\code{PEFT/QLoRA} (LLM). Suivi des expériences :
\code{matplotlib} pour les courbes, export des métriques en CSV.

% =====================================================================
\section{Le mode \og réseaux sociaux \fg{} : réalité et stratégie}

\begin{alerte}
En 2026, le \textbf{scraping de Facebook déconnecté est bloqué} : \code{mbasic}
renvoie une erreur HTTP~400. Toute collecte réelle exige une \textbf{session
authentifiée} (donc hors CGU), ou une \textbf{API tierce payante}.
\end{alerte}

Stratégie pragmatique pour ne \textbf{jamais} dépendre du scraping le jour de la
soutenance :
\begin{enumerate}[leftmargin=1.2cm,nosep]
  \item \textbf{Dataset de secours} (recommandé) : un corpus de tweets / avis FR
        déjà collecté, chargé dans le mode \og réseaux sociaux \fg{}. La démo
        fonctionne à coup sûr.
  \item \textbf{Scraper (apprentissage)} : module \code{scraper\_fb.py} déjà en
        place (\code{sentiment/}). Fonctionne avec un \textbf{compte dédié} +
        cookie de session, ou via \textbf{Playwright} pour le contenu dynamique.
        À présenter avec ses limites et le cadrage \textbf{RGPD} (Pages publiques
        seulement, anonymisation par hachage, faible débit).
  \item \textbf{Option robuste} : API tierce (Apify) si un petit budget existe.
\end{enumerate}

% =====================================================================
\section{Pile technique et organisation du projet}

\subsection{Pile}
Python 3.12 (\code{uv}) ; \code{scikit-learn}, \code{pandas}, \code{spaCy}/\code{nltk} ;
\code{transformers}~+~\code{torch} ; \code{unsloth} (option) ; \code{streamlit} ;
\code{matplotlib}/\code{plotly}, \code{wordcloud} ; \code{requests}/
\code{beautifulsoup4}/\code{playwright} (collecte).

\subsection{Arborescence cible}
{\scriptsize
\begin{verbatim}
datamining/sentiment/
  app.py                  # interface Streamlit (3 onglets)
  pretraitement.py        # nettoyage FR, negation, emojis
  moteur_lexical.py       # approche lexicale (FEEL)
  moteur_ml.py            # TF-IDF + classifieurs sklearn
  moteur_transformer.py   # CamemBERT (Hugging Face)
  entrainement/           # scripts d'entrainement + benchmark
  scraper_fb.py           # collecte reseaux sociaux (deja present)
  data/                   # datasets FR + corpus maison annote
  modeles/                # modeles entraines sauvegardes
  resultats/              # metriques, matrices de confusion
\end{verbatim}}

% =====================================================================
\section{Plan de réalisation (jalons)}

\begin{center}
\small
\begin{tabularx}{\linewidth}{c >{\bfseries}l X}
\toprule
\# & Jalon & Contenu \\
\midrule
1 & Prétraitement + lexique & Baseline explicable qui tourne de bout en bout \\
2 & ML classique & TF-IDF + classifieurs, dataset FR, premier benchmark \\
3 & CamemBERT & Fine-tuning Hugging Face, comparaison \\
4 & (option) LLM + Unsloth & QLoRA, justification + aspects \\
5 & Interface Streamlit & Les 3 modes + sélecteur de moteur \\
6 & Mode réseaux sociaux & Dataset de secours + scraper + RGPD \\
7 & Évaluation + rapport & Modèle optimal, visualisations, rédaction \\
\bottomrule
\end{tabularx}
\end{center}

% =====================================================================
\section{Risques et limites à documenter}
\begin{itemize}[leftmargin=1.2cm,nosep]
  \item \textbf{Ironie / sarcasme} : difficiles pour tous les modèles ; à mesurer.
  \item \textbf{Ressources FR} limitées vs anglais : lexiques et corpus plus rares.
  \item \textbf{RGPD} : données personnelles $\rightarrow$ anonymisation, base légale,
        Pages publiques uniquement.
  \item \textbf{Dépendance au scraping} : neutralisée par le dataset de secours.
  \item \textbf{Échantillon non représentatif} : à nuancer dans l'interprétation.
\end{itemize}

% =====================================================================
\section{Synthèse des décisions}
\begin{enumerate}[leftmargin=1.2cm,nosep]
  \item Système : \textbf{analyse de sentiments FR}, application \textbf{Streamlit}.
  \item \textbf{3 modes d'entrée} (phrase / texte / réseaux sociaux) conformes au sujet.
  \item \textbf{2 approches} implémentées et comparées : \textbf{lexicale} et \textbf{ML}
        (classique + CamemBERT), avec option \textbf{LLM + Unsloth}.
  \item Modèle \textbf{optimal} choisi par \textbf{benchmark chiffré} (F1-macro).
  \item Mode réseaux sociaux sécurisé par un \textbf{dataset de secours} + scraper cadré RGPD.
  \item Finalité : \textbf{aide à la décision} (visualisations, export).
\end{enumerate}

\vfill
\begin{center}\color{gris}\small
Document de brainstorming -- base de discussion pour la conception détaillée.
\end{center}

\end{document}
